% ---------- Datei: chapters/02_analyse.tex ----------
\chapter{Analysephase}

\section{Ist-Zustand und Problemanalyse}
Die Söldner Consult GmbH ist ein IT-Dienstleistungsunternehmen mit Schwerpunkt auf Cloud-Architekturen, Netzwerkdesign und Infrastrukturmodernisierung. Der Kunde ist ein mittelständisches Unternehmen mit Hauptstandort in Nürnberg sowie zwei Außenstellen.

Die bestehende Infrastruktur umfasst:
\begin{itemize}
    \item ein lokales Rechenzentrum,
    \item mehrere VPN-Anbindungen,
    \item sowie Cloud-Ressourcen innerhalb der Google Cloud Platform.
\end{itemize}

Die aktuelle Netzwerkstruktur basiert auf mehreren separaten Site-to-Site-VPN-Verbindungen. Das Unternehmen plant eine schrittweise Migration weiterer Dienste in die Cloud, jedoch erfüllt die bestehende Netzwerkarchitektur die Anforderungen an Skalierbarkeit und zentrale Verwaltung nicht mehr.

Probleme der bestehenden Architektur:
\begin{itemize}
    \item manuelle Pflege statischer Routen,
    \item fehlende zentrale Verwaltung,
    \item hohe Komplexität bei Erweiterungen,
    \item eingeschränkte Transparenz,
    \item inkonsistente Sicherheitsrichtlinien.
\end{itemize}

Zusätzlich existieren keine zentralen Monitoring-Mechanismen zur Überwachung der hybriden Netzwerkverbindungen. Die bestehende Architektur verursacht erhöhten administrativen Aufwand und erschwert zukünftige Cloud-Erweiterungen.

\section{Problemstellung}
Mit zunehmender Anzahl angebundener Standorte steigt die Komplexität der bestehenden VPN-Mesh-Struktur erheblich.

Neue Standorte erfordern:
\begin{itemize}
    \item zusätzliche Tunnel,
    \item manuelle Routing-Anpassungen,
    \item individuelle Firewall-Konfigurationen.
\end{itemize}

Dies erhöht die Fehleranfälligkeit, die Bereitstellungszeit sowie den Betriebsaufwand. Zudem bestehen Risiken hinsichtlich fehlender Redundanz, mangelnder Übersichtlichkeit und eingeschränkter Skalierbarkeit.

\section{Projektziel (Soll-Konzept)}
Ziel des Projekts ist die Implementierung einer zentralen hybriden Netzwerkarchitektur innerhalb der Google Cloud.

Hierfür soll:
\begin{itemize}
    \item ein NCC-Hub eingerichtet,
    \item zwei VPC-Netzwerke angebunden,
    \item zwei On-Premises-Standorte per HA-VPN integriert,
    \item dynamisches Routing mittels BGP implementiert,
    \item ein zentrales Sicherheitskonzept umgesetzt,
    \item sowie Monitoring und Logging eingerichtet werden.
\end{itemize}

Die Zielarchitektur soll hochverfügbar, skalierbar, standardisiert sowie einfach erweiterbar sein.

\section{Anforderungsanalyse}

\textbf{Funktionale Anforderungen:}
\begin{itemize}
    \item Kommunikation zwischen allen Standorten
    \item Zentrale Verwaltung der Netzwerkverbindungen
    \item Dynamisches Routing
    \item Hochverfügbarkeit der VPN-Verbindungen
    \item Zentrale Sicherheitsrichtlinien
    \item Monitoring und Logging
\end{itemize}

\textbf{Nichtfunktionale Anforderungen:}
\begin{itemize}
    \item Skalierbarkeit
    \item Wartbarkeit
    \item Ausfallsicherheit
    \item geringe Administrationskomplexität
    \item DSGVO-konforme Protokollierung
\end{itemize}

\subsection{Lasten- und Pflichtenheft}
\begin{table}[!ht]
\centering
\begin{tabular}{|p{0.25\textwidth}|p{0.3\textwidth}|p{0.35\textwidth}|}
\hline
\textbf{Kategorie} & \textbf{Lastenheft (Kunde)} & \textbf{Pflichtenheft (Lösung)} \\ \hline
Architektur & Zentrale Verwaltung & Aufbau eines GCP NCC Hubs \\ \hline
Verfügbarkeit & Ausfallsichere Anbindung & Einsatz von HA Cloud VPN \\ \hline
Routing & Automatischer Routenaustausch & Cloud Router mit BGP \\ \hline
Sicherheit & Zentraler Schutz & GCP Firewall-Regeln (Default Deny) \\ \hline
\end{tabular}
\caption{Tabellarisches Lasten- und Pflichtenheft}
\end{table}

\section{Stakeholder-Analyse}
\begin{table}[!ht]
\centering
\begin{tabular}{|p{0.2\textwidth}|p{0.15\textwidth}|l|p{0.35\textwidth}|}
\hline
\textbf{Stakeholder} & \textbf{Rolle} & \textbf{Einfluss} & \textbf{Interesse} \\ \hline
N. Kunak & Umsetzung & Hoch & Projekterfolg, IHK-Abschluss \\ \hline
C. Söldner & Betreuer & Hoch & Qualität, Kundenzufriedenheit \\ \hline
Kunde (IT) & Auftraggeber & Hoch & Stabile, sichere Infrastruktur \\ \hline
Netzwerk-Admins & Anwender & Mittel & Geringer Wartungsaufwand \\ \hline
\end{tabular}
\caption{Stakeholder-Analyse}
\end{table}

\section{Wirtschaftlichkeitsbetrachtung}
Die bisherige Architektur verursacht durch die manuelle Einrichtung von VPN-Tunneln und statischen Routen einen hohen administrativen Aufwand. Die neue NCC-basierte Architektur reduziert Fehlerquellen und verkürzt die Bereitstellungszeiten für neue Standorte erheblich.

Die internen Personalkosten für die Projektumsetzung basieren auf einem fiktiven, marktüblichen Stundensatz für Systemadministratoren in Höhe von 70\,€.

\textbf{Kostenübersicht (Einmalkosten):}
\begin{table}[!ht]
\centering
\begin{tabular}{|p{0.6\textwidth}|r|}
\hline
\textbf{Position} & \textbf{Kosten} \\ \hline
Projektarbeitszeit (40 Stunden $\times$ 70\,€) & 2.800,00\,€ \\ \hline
Google Cloud Ressourcen (Einrichtung NCC Hub, Router) & 320,00\,€ \\ \hline
Testumgebung (Temporäre VMs und Traffic) & 180,00\,€ \\ \hline
\textbf{Gesamte Projektkosten} & \textbf{3.300,00\,€} \\ \hline
\end{tabular}
\caption{Einmalige Projektkosten}
\end{table}

\textbf{Laufende Betriebskosten (Google Cloud) pro Jahr:}
\begin{table}[!ht]
\centering
\begin{tabular}{|p{0.6\textwidth}|r|}
\hline
\textbf{Position} & \textbf{Kosten} \\ \hline
NCC Hub \& Spoke-Verbindungen (Grundgebühr) & 1.200,00\,€ \\ \hline
HA VPN Tunnel (2 Standorte, redundante Tunnel) & 1.400,00\,€ \\ \hline
Traffic (Egress-Datenverkehr, geschätzt) & 600,00\,€ \\ \hline
\textbf{Gesamte Betriebskosten (p.\,a.)} & \textbf{3.200,00\,€} \\ \hline
\end{tabular}
\caption{Laufende Cloud-Kosten pro Jahr}
\end{table}

\textbf{Erwartete Einsparungen pro Jahr (Amortisation):} \\
Durch die Automatisierung des Routings und die vereinfachte Verwaltung sinkt der monatliche Wartungsaufwand um schätzungsweise 6 Stunden. Zudem können geplante Standortanbindungen schneller realisiert werden.

\begin{table}[!ht]
\centering
\begin{tabular}{|p{0.6\textwidth}|r|}
\hline
\textbf{Bereich} & \textbf{Einsparung} \\ \hline
Reduzierter Administrationsaufwand (72 Stunden $\times$ 70\,€) & 5.040,00\,€ \\ \hline
Vermeidung von Ausfallzeiten durch HA VPN (geschätzt) & 1.200,00\,€ \\ \hline
\textbf{Gesamte Einsparungen (p.\,a.)} & \textbf{6.240,00\,€} \\ \hline
\end{tabular}
\caption{Erwartete jährliche Einsparungen}
\end{table}

\textbf{Fazit der Wirtschaftlichkeit:} \\
Den laufenden Cloud-Kosten von 3.200,00\,€ stehen jährliche Einsparungen in Höhe von 6.240,00\,€ gegenüber. Es ergibt sich ein positiver Nettonutzen von \textbf{3.040,00\,€ pro Jahr}. Die initialen Projektkosten von 3.300,00\,€ amortisieren sich somit bereits nach \textbf{knapp 13 Monaten} (Return on Investment). Das Projekt ist aus wirtschaftlicher Sicht sehr profitabel.

\section{Nutzwertanalyse}
\begin{table}[!ht]
\centering
\begin{tabular}{|p{0.3\textwidth}|c|c|c|}
\hline
\textbf{Kriterium} & \textbf{Gewichtung} & \textbf{Klassische VPN-Mesh} & \textbf{NCC-Architektur} \\ \hline
Skalierbarkeit & 30\,\% & 2 & 5 \\ \hline
Wartbarkeit & 25\,\% & 2 & 5 \\ \hline
Hochverfügbarkeit & 20\,\% & 3 & 5 \\ \hline
Administrationsaufwand & 15\,\% & 2 & 5 \\ \hline
Transparenz & 10\,\% & 2 & 4 \\ \hline
\textbf{Gesamtnutzen} & \textbf{100\,\%} & \textbf{2,15} & \textbf{4,9} \\ \hline
\end{tabular}
\caption{Entscheidungsmatrix}
\end{table}
Die NCC-Architektur erzielt den deutlich höheren Gesamtnutzen.

\section{Risikoanalyse}
Während der Projektumsetzung können Risiken auftreten. Diese wurden im Vorfeld bewertet, um Gegenmaßnahmen zu definieren. Die Bewertung erfolgt auf einer Skala von 1 (niedrig) bis 3 (hoch).

\begin{table}[htbp]
\centering
\begin{tabular}{|p{0.35\textwidth}|p{0.1\textwidth}|p{0.15\textwidth}|p{0.3\textwidth}|}
\hline
\textbf{Risiko} & \textbf{Eintritt} & \textbf{Auswirkung} & \textbf{Gegenmaßnahme} \\ \hline
\raggedright Ausfall bestehender VPN-Tunnel bei der Migration & 2 & 3 & \raggedright Durchführung der Umstellung in einem definierten Wartungsfenster am Wochenende. \tabularnewline \hline
\raggedright Fehlerhafte BGP-Routen blockieren den Traffic & 2 & 3 & \raggedright Vorheriger Test in einer isolierten Sandbox-Umgebung mit Test-VMs. \tabularnewline \hline
\raggedright Zeitüberschreitung (40h Limit) & 1 & 2 & \raggedright Strikte Einhaltung des Projektstrukturplans und Einplanung von Pufferzeiten. \tabularnewline \hline
\end{tabular}
\caption{Risikoanalyse}
\end{table}

